% part: intuitionistic-logic
% chapter: introduction
% section: syntax

\documentclass[../../../include/open-logic-section]{subfiles}

\begin{document}

\olfileid{int}{int}{syn}

\olsection{સ્ફુરણાવાદી તર્કની વાક્યરચના}

સ્ફુરણાવાદી તર્કની વાક્યરચના વિધાનાત્મક તર્ક જેવી જ છે.
પ્રશિષ્ટ વિધાનાત્મક તર્કમાં કેટલાક સંયોજકોને બીજાં સંયોજકો
વડે વ્યાખ્યાયિત કરી શકાય છે. દાખલા તરીકે, $!A \lif !B$ને
$\lnot !A \lor !B$ વડે, અથવા $!A \lor !B$ને
$\lnot(\lnot !A \land \lnot !B)$ વડે વ્યાખ્યાયિત કરી શકીએ.
તેથી પ્રશિષ્ટ તર્કની રજૂઆતોમાં કેટલાક સંયોજકો આવી વ્યાખ્યાઓના
સંક્ષેપ તરીકે રજૂ થાય છે. સ્ફુરણાવાદી તર્કમાં એવું નથી, બે
અપવાદો સિવાય: $\lnot !A$ને $!A \lif \lfalse$ના સંક્ષેપ
તરીકે વ્યાખ્યાયિત કરી શકાય છે---અને ઘણી વાર કરાય છે. તો,
અલબત્ત, $\lfalse$ પોતે વ્યાખ્યાયિત સંક્ષેપ ન હોવું જોઈએ!{}
વળી, પ્રશિષ્ટ તર્કની જેમ $!A \liff !B$ને
$(!A \lif !B) \land (!B \lif !A)$ તરીકે વ્યાખ્યાયિત કરી
શકાય છે.

વિધાનાત્મક સ્ફુરણાવાદી તર્કનાં !!^{formula}s
\emph{!!{propositional variable}s} અને વિધાનાત્મક
અચલાંક~$\lfalse$માંથી \emph{તાર્કિક સંયોજકો} વડે રચાય છે.
આપણી પાસે છે:

\begin{enumerate}
\item !!{propositional variable}sનો !!^a{denumerable} ગણ~$\PVar$:
  $\Obj p_0$,
  $\Obj p_1$, \dots
  \item !!{falsity} માટેનો વિધાનાત્મક અચલાંક~$\lfalse$.
\item તાર્કિક સંયોજકો: $\land$ (સંયોજન), $\lor$ (વિયોજન),
  $\lif$ (શરતી સંયોજક).
\item વિરામચિહ્નો: (, ), અને અલ્પવિરામ.
\end{enumerate}

\begin{defn}[સૂત્ર]
\ollabel{defn:formulas} વિધાનાત્મક સ્ફુરણાવાદી તર્કનાં
\emph{!!{formula}s}નો ગણ~$\Frm[L_0]$ નીચે મુજબ આગમનાત્મક
રીતે વ્યાખ્યાયિત થાય છે:
\begin{enumerate}
\item $\lfalse$ પરમાણુક !!{formula} છે.

\item દરેક !!{propositional variable}~$\Obj p_i$ પરમાણુક
  !!{formula} છે.

\item જો $!A$ અને $!B$ એ !!{formula}s હોય, તો $(!A \land
  !B)$ પણ !!a{formula} છે.

\item જો $!A$ અને $!B$ એ !!{formula}s હોય, તો $(!A \lor !B)$
  પણ !!a{formula} છે.

\item જો $!A$ અને $!B$ એ !!{formula}s હોય, તો $(!A \lif !B)$
  પણ !!a{formula} છે.

\tagitem{limitClause}{આ સિવાય બીજું કશું !!a{formula} નથી.}{}
\end{enumerate}
\end{defn}

ઉપર જણાવેલા મૂળભૂત સંયોજકો ઉપરાંત, નીચેના \emph{વ્યાખ્યાયિત}
સંકેતો પણ વાપરીએ છીએ: $\lnot$ (નિષેધ) અને $\liff$
(!!{biconditional}). આ વ્યાખ્યાયિત સંક્રિયકો વડે બનેલાં
સૂત્રોને નીચે મુજબ સમજવાનાં છે:
\begin{enumerate}
\item $\lnot !A$ એ $!A \lif \lfalse$નો સંક્ષેપ છે.

\item $!A \liff !B$ એ $(!A \lif !B) \land (!B
  \lif !A)$નો સંક્ષેપ છે.
\end{enumerate}

જોકે $\lnot$ને ઔપચારિક રીતે સંક્ષેપ ગણીએ છીએ, કેટલીક વાર
તે મૂળભૂત હોય તેમ $\lnot$ માટે સ્પષ્ટ નિયમો અને વ્યાખ્યાની
શરતો આપીશું. મુખ્યત્વે અભ્યાસ-પ્રશ્નો રજૂ કરી શકાય તે માટે
આવું કરીશું.

\end{document}
